\documentclass[aspectratio=169,xcolor=dvipsnames, t]{beamer}
\usepackage{fontspec}
\usepackage{beamerthemeSimplePlusAIC}
\usepackage[utf8]{inputenc}
\usepackage[T1]{fontenc}
\usepackage[spanish]{babel}
\usepackage{hyperref}
\usepackage{listings}
\usepackage{xcolor}
\usepackage{polyglossia}
\usepackage{graphicx}
\usepackage{booktabs}
\usepackage{svg}
\usepackage{tikz}
\usepackage{makecell}
\usepackage{wrapfig}

\title[Agentes]{Sistemas Basados en Agentes}
\subtitle{Arquitectura y Patrones}
\author[Jorge]{Inteligencia Artificial: Dr. Jorge Roa | Dra. Milagros Gutiérrez}
\institute[CIDISI]{CIDISI - UTN - FRSF}
\date{Mayo 2025}

\begin{document}
\maketitlepage

%%%%%%%%%%%%%%%%%%%%%%%%
\begin{frame}
\frametitle{Agenda}
\tableofcontents
\end{frame}

\section{Introducción y Motivación}


%%%%%%%%%%%%%%%%%%%%%%%%
\begin{frame}
\frametitle{¿Qué es un agente LLM?}
\begin{itemize}
    \item Es un sistema autónomo basado en un modelo de lenguaje.
    \item Va más allá de generar texto: puede \textbf{razonar, planear, actuar y adaptarse}.
    \item Integra múltiples componentes: LLM, herramientas, memoria, planificación, etc.
    \item Ejecuta tareas complejas en ciclos iterativos hasta cumplir un objetivo.
\end{itemize}
\end{frame}

%%%%%%%%%%%%%%%%%%%%%%%%
\begin{frame}
\frametitle{¿Por qué agentes basados en LLM?}
\begin{itemize}
    \item Los LLMs han revolucionado la IA, pero los agentes LLM llevan la autonomía al siguiente nivel.
    \item Permiten resolver tareas complejas, multi-paso y adaptativas, integrando memoria, razonamiento y herramientas.
    \item Transforman industrias: salud, finanzas, logística, educación, software, etc.
\end{itemize}
\end{frame}

%%%%%%%%%%%%%%%%%%%%%%%%
\begin{frame}
\frametitle{De LLMs básicos a agentes autónomos}
\begin{itemize}
    \item LLM básico: responde a un prompt, sin memoria ni acción.
    \item Agente LLM: descompone tareas, usa herramientas, mantiene contexto, se auto-optimiza.
    \item Ejemplo: De un chatbot a un asistente que agenda, consulta bases y ejecuta acciones.
\end{itemize}
\end{frame}

\section{Visión General de la Arquitectura de Agentes LLM}

%%%%%%%%%%%%%%%%%%%%%%%%
\begin{frame}
\frametitle{Arquitectura general de un agente LLM}
\begin{itemize}
    \item Los agentes LLM son sistemas modulares.
    \item Integran varios componentes centrales para lograr autonomía y adaptabilidad.
    \item Cada componente cumple un rol específico en el ciclo de percepción, razonamiento y acción.
\end{itemize}
\begin{figure}[ht]
    \centering
    \includegraphics[width=0.62\textwidth]{figures/coreAI.jpg}
\end{figure}

\end{frame}


% \section{Componentes principales de un agente}

%%%%%%%%%%%%%%%%%%%%%%%%
\begin{frame}
\frametitle{Componentes principales de un agente}
\begin{enumerate}
    \item \textbf{LLM}: núcleo cognitivo que interpreta, razona y genera.
    \item \textbf{Memoria}: retiene contexto (corto y largo plazo).
    \item \textbf{Retriever}: accede a fuentes externas de conocimiento.
    \item \textbf{Herramientas}: ejecutan acciones en el mundo real (API, código, búsqueda).
    \item \textbf{Planificador}: define el orden de pasos.
    \item \textbf{Loop de reflexión}: mejora resultados iterativamente.
    \item \textbf{Seguridad y control}: limita acciones peligrosas o erróneas.
\end{enumerate}
\end{frame}

%%%%%%%%%%%%%%%%%%%%%%%%
\begin{frame}
\frametitle{1. LLM como cerebro del agente}
\begin{itemize}
    \item El LLM (ej: GPT-4, Claude, LLaMA) interpreta el input del usuario.
    \item Formula pensamientos intermedios y decide qué hacer a continuación.
    \item Puede generar instrucciones para herramientas u otros módulos.
    \item Todo el sistema gira en torno a su capacidad de razonamiento.
\end{itemize}
\end{frame}

%%%%%%%%%%%%%%%%%%%%%%%%
\begin{frame}
\frametitle{2. Memoria: corto y largo plazo}
\begin{itemize}
    \item \textbf{Corto plazo}: guarda información del ciclo actual o sesión.
    \item \textbf{Largo plazo}: almacena preferencias, historial, aprendizajes.
    \item Mejora la personalización y continuidad.
    \item Puede almacenarse como texto, embeddings, logs estructurados.
\end{itemize}
\end{frame}

%%%%%%%%%%%%%%%%%%%%%%%%
\begin{frame}
\frametitle{3. Retriever y base de conocimiento}
\begin{itemize}
    \item Extrae datos relevantes desde bases vectoriales, bases SQL, la web o documentos.
    \item Apoya la generación con \textbf{información externa actualizada}.
    \item Es fundamental para agentes que usan RAG (retrieval augmented generation).
\end{itemize}
\end{frame}

%%%%%%%%%%%%%%%%%%%%%%%%
\begin{frame}
\frametitle{4. Herramientas y acciones}
\begin{columns}[c] % top-aligned columns
    \begin{column}{0.6\textwidth}
        \begin{itemize}
            \item Permiten que el agente interactúe con el entorno.
            \item Tipos: APIs REST, ejecución de código, comandos, navegadores, bases de datos.
            \item Decidir cuándo y cómo usarlas es una parte crítica del razonamiento del agente.
        \end{itemize}
    \end{column}
    \begin{column}{0.4\textwidth}
        \centering
        \includegraphics[width=\textwidth]{figures/tool_pattern.png}
    \end{column}
\end{columns}
\end{frame}

%%%%%%%%%%%%%%%%%%%%%%%%
\begin{frame}
\frametitle{5. Planificador}
\begin{columns}[c] % top-aligned columns
    \begin{column}{0.6\textwidth}
        \begin{itemize}
            \item Descompone tareas complejas en subtareas.
            \item Puede usar planificación jerárquica, "chain-of-thought", o árboles de decisiones.
            \item Determina el orden, dependencias, y cuándo paralelizar tareas.
        \end{itemize}
    \end{column}
    \begin{column}{0.4\textwidth}
        \centering
        \includegraphics[width=\textwidth]{figures/planning_pattern.png}
    \end{column}
\end{columns}
\end{frame}

%%%%%%%%%%%%%%%%%%%%%%%%
\begin{frame}
\frametitle{6. Loop de reflexión y feedback}
\begin{columns}[c] % align text and image at the top
    \begin{column}{0.6\textwidth}
        \begin{itemize}
            \item El agente analiza resultados y decide si continuar, corregir o replantear.
            \item Ciclo típico: \textbf{Thought → Action → Observation}.
            \item Permite corrección de errores, adaptación dinámica y mejora continua.
        \end{itemize}
    \end{column}
    \begin{column}{0.4\textwidth}
        \centering
        \includegraphics[width=\textwidth]{figures/reflection_pattern.png}
    \end{column}
\end{columns}
\end{frame}

%%%%%%%%%%%%%%%%%%%%%%%%
\begin{frame}
\frametitle{7. Módulo de seguridad y control}
\begin{itemize}
    \item Verifica que las acciones del agente sean seguras, autorizadas y adecuadas.
    \item Implementa filtros, validaciones, límites de permisos o revisiones humanas.
    \item Esencial en entornos sensibles: salud, finanzas, software crítico.
\end{itemize}
\end{frame}



\section{Ciclo de Vida y Workflow de un Agente LLM}

%%%%%%%%%%%%%%%%%%%%%%%%
\begin{frame}
\frametitle{Ciclo de vida de un agente LLM}
\begin{enumerate}
    \item Inicialización: prompt del sistema y análisis de la consulta.
    \item Planificación: generación de plan multi-paso.
    \item Acción: ejecución de herramientas/APIs.
    \item Observación: análisis de resultados y actualización de contexto.
    \item Re-planificación: ajuste de estrategia según feedback.
\end{enumerate}
\end{frame}

%%%%%%%%%%%%%%%%%%%%%%%%
\begin{frame}
\frametitle{Ciclo iterativo Thought-Action-Observation}
\begin{itemize}
    \item Thought: razonamiento y decisión.
    \item Action: ejecución de acción/herramienta.
    \item Observation: análisis de resultados y actualización de memoria.
    \item El ciclo se repite hasta cumplir el objetivo.
\end{itemize}
\end{frame}

\section{Patrones de Agentes LLM}

%%%%%%%%%%%%%%%%%%%%%%%%
\begin{frame}
\frametitle{Patrones de agentes: visión general}
\begin{itemize}
    \item Los patrones definen cómo se estructuran y coordinan los agentes.
    \item Permiten modularidad, especialización, escalabilidad y robustez.
    \item Se pueden combinar para sistemas complejos.
\end{itemize}
\end{frame}

%%%%%%%%%%%%%%%%%%%%%%%%
\begin{frame}
\frametitle{Patrón: Prompt Chaining}
\begin{itemize}
    \item Encadena prompts para descomponer tareas complejas.
    \item Cada paso valida y refina el resultado anterior.
    \item Ejemplo: generación, traducción y revisión de un documento.
\end{itemize}
\end{frame}

%%%%%%%%%%%%%%%%%%%%%%%%
\begin{frame}
\frametitle{Patrón: Router Architecture}
\begin{itemize}
    \item El agente decide a qué herramienta, skill o sub-agente enviar cada input.
    \item Usa salidas estructuradas para garantizar interpretabilidad.
    \item Ejemplo: elegir entre APIs de clima, finanzas o calendario.
\end{itemize}
\end{frame}

%%%%%%%%%%%%%%%%%%%%%%%%
\begin{frame}
\frametitle{Patrón: Orchestrator Architecture}
\begin{itemize}
    \item Coordina múltiples herramientas o skills en secuencia planificada.
    \item Maneja dependencias y resultados intermedios.
    \item Ejemplo: pipeline de análisis de datos con varios pasos.
\end{itemize}
\begin{figure}[ht]
    \centering
    \includegraphics[width=0.8\textwidth]{figures/workflow.png}
\end{figure}

\end{frame}

%%%%%%%%%%%%%%%%%%%%%%%%
\begin{frame}
\frametitle{Patrón: Parallel Execution}
\begin{itemize}
    \item Ejecuta subtareas independientes en paralelo.
    \item Aumenta velocidad y throughput.
    \item Ejemplo: procesamiento batch de consultas o soporte multi-canal.
\end{itemize}

\end{frame}

%%%%%%%%%%%%%%%%%%%%%%%%
\begin{frame}
\frametitle{Patrón: Multi-Agent Pattern}

\begin{columns}[c] % [c] para centrar verticalmente
    \begin{column}{0.5\textwidth}
        \begin{itemize}
            \item Varios agentes especializados colaboran para un objetivo común.
            \item Puede ser centralizado (manager) o descentralizado.
            \item Ejemplo: Project Manager, Coder, Tester y Critic trabajando juntos.
        \end{itemize}
    \end{column}

    \begin{column}{0.4\textwidth}
        \includegraphics[width=\textwidth]{figures/multiagent_pattern.png}
    \end{column}
\end{columns}

\end{frame}

\begin{frame}
\frametitle{Patrón: Structured Output \& Tool Calling}

\begin{columns}[c]  % [c] centra verticalmente el contenido
    \begin{column}{0.5\textwidth}
        \begin{itemize}
            \item El agente produce salidas en formatos predefinidos.
            \item Permite invocación confiable de herramientas y APIs.
            \item Ejemplo: output JSON para llamar a una función externa.
        \end{itemize}
    \end{column}

    \begin{column}{0.4\textwidth}
        \includegraphics[width=\textwidth]{figures/tool_pattern.png}
    \end{column}
\end{columns}

\end{frame}

%%%%%%%%%%%%%%%%%%%%%%%%
\begin{frame}
\frametitle{Patrón: User-in-the-Loop}
\begin{itemize}
    \item Integra feedback humano durante o después de la ejecución.
    \item Balancea autonomía y control.
    \item Ejemplo: revisión humana de decisiones críticas.
\end{itemize}
\end{frame}

%%%%%%%%%%%%%%%%%%%%%%%%
\begin{frame}
\frametitle{Patrón: Memory-Augmented Agents}
\begin{itemize}
    \item Mantienen memoria persistente entre sesiones.
    \item Mejoran coherencia y personalización.
    \item Ejemplo: bot de soporte que recuerda interacciones previas.
\end{itemize}
\end{frame}

\section{Memoria y Personalización}

%%%%%%%%%%%%%%%%%%%%%%%%
\begin{frame}
\frametitle{Memoria a corto y largo plazo}
\begin{itemize}
    \item \textbf{Corto plazo}: contexto inmediato, relevante para la sesión actual.
    \item \textbf{Largo plazo}: preferencias, historial, conocimiento duradero.
    \item Tipos adicionales: episódica, semántica, procedimental, de entidades.
\end{itemize}
\end{frame}

%%%%%%%%%%%%%%%%%%%%%%%%
\begin{frame}
\frametitle{Persistencia de memoria y personalización}
\begin{itemize}
    \item Permite continuidad de contexto y relaciones genuinas.
    \item Recomendaciones personalizadas y reducción de carga cognitiva.
    \item Soporte para workflows multi-sesión y reflexión.
    \item Evita inconsistencias y mejora la credibilidad.
\end{itemize}
\end{frame}

\section{Aplicaciones y Futuro}

%%%%%%%%%%%%%%%%%%%%%%%%
\begin{frame}
\frametitle{Aplicaciones reales de agentes LLM}
\begin{itemize}
    \item Salud: Dr. Chatbot, síntesis de historias clínicas.
    \item Finanzas: análisis de mercado, monitoreo de sentimiento.
    \item Manufactura: troubleshooting, generación de código.
    \item Logística: optimización de rutas, coordinación en tiempo real.
    \item Retail: soporte al cliente, marketing personalizado.
    \item Capacitación: aprendizaje adaptativo.
\end{itemize}
\end{frame}

%%%%%%%%%%%%%%%%%%%%%%%%
\begin{frame}
\frametitle{Desafíos y futuro de los agentes LLM}
\begin{itemize}
    \item Gestión de contexto y memoria a gran escala.
    \item Seguridad, ética y control humano.
    \item Evaluación robusta y orquestación multi-agente.
    \item Integración multimodal y avance hacia AGI.
\end{itemize}
\end{frame}

%%%%%%%%%%%%%%%%%%%%%%%%
\begin{frame}
\frametitle{Conclusiones}
\begin{itemize}
    \item Los agentes LLM son el futuro de la automatización inteligente.
    \item Integran razonamiento, memoria, herramientas y colaboración multi-agente.
    \item Permiten resolver tareas complejas, personalizar experiencias y escalar soluciones en el mundo real.
\end{itemize}
\end{frame}

\end{document}